\label{sec:well-ordering-principle}

\begin{topicbox}{The Well-Ordering Principle}
\begin{exposition}
The well-ordering principle asserts that every nonempty subset of $\mathbb{N}$
contains a least element. This is not merely a consequence of the order axioms;
it is a property that distinguishes $\mathbb{N}$ from dense ordered sets such
as $\mathbb{Q}$, which have no smallest positive element. The well-ordering
principle is equivalent to ordinary induction over $\mathbb{N}$, but it
supports a distinct proof style: to establish a universal statement, assume
it fails for some element and derive a contradiction by examining the least
failure.
\end{exposition}

\begin{theorembox}{Theorem (Well-Ordering Principle)}
\begin{theorem}[Well-Ordering Principle]
\label{thm:well-ordering-principle}
Let $(P,S,1)$ be a Peano system. Every nonempty subset of $P$ has a least
element with respect to the order $\le$ on $P$.
\hyperref[prf:well-ordering-principle]{Go to proof.}
\end{theorem}
\end{theorembox}

\begin{remark*}[Standard quantified statement]
\[
\begin{aligned}
&\text{Let } (P,S,1) \text{ be a Peano system.}\\
&\forall A \subseteq P\;
\Bigl(A \neq \varnothing \Longrightarrow
\exists m \in A\;\forall a \in A\;(m \le a)\Bigr).
\end{aligned}
\]
\end{remark*}

\begin{remark*}[Predicate reading]
\[
\operatorname{WellOrderingPrinciple}
\;\coloneqq\;
\forall A\;\bigl((A \subseteq P \land \exists a\,(a \in A))
\Rightarrow \exists m\,(m \in A \land \forall a\,(a \in A \Rightarrow m \le a))\bigr).
\]
\end{remark*}

\begin{remark*}[Negated quantified statement]
\[
\begin{aligned}
&\text{Let } (P,S,1) \text{ be a Peano system.}\\
&\exists A \subseteq P\;
\Bigl(A \neq \varnothing \land
\forall m \in A\;\exists a \in A\;(a < m)\Bigr).
\end{aligned}
\]
\end{remark*}

\begin{remark*}[Failure modes]
The well-ordering principle would fail if some nonempty subset of $P$ had no
least element. That means for every candidate minimum, a strictly smaller
element of the subset exists, producing an infinite strictly descending chain
in $P$. Induction rules this out.
\end{remark*}

\begin{remark*}[Failure mode decomposition]
\[
\neg\operatorname{WellOrderingPrinciple}
\Longleftrightarrow
\exists A \subseteq P\;
\Bigl(\underbrace{A \neq \varnothing}_{\text{nonempty}}
\land
\underbrace{\forall m \in A\;\exists a \in A\;(a < m)}_{\text{no least element}}
\Bigr).
\]
\end{remark*}

\begin{remark*}[Interpretation]
The well-ordering principle is the order-theoretic complement to induction.
Induction says one can build up: start at $1$ and keep stepping forward.
Well-ordering says one can always descend to a minimum: any nonempty set has
a floor. Proof strategies that use well-ordering typically assume a property
fails somewhere, form the set of failures, extract its minimum by
well-ordering, and derive a contradiction from the minimality of that element.
\end{remark*}

\begin{dependencies}
\begin{itemize}
  \item \hyperref[def:peano-system]{Peano System}
  \item \hyperref[def:less-than-or-equal-on-peano-system]{Less Than or Equal on a Peano System}
  \item \hyperref[thm:induction-principle-for-peano-system]{Induction Principle for a Peano System}
  \item \hyperref[thm:trichotomy-on-peano-system]{Trichotomy on a Peano System}
\end{itemize}
\end{dependencies}
\end{topicbox}

\begin{topicbox}{Equivalence of Induction, Strong Induction, and Well-Ordering}
\begin{exposition}
Ordinary induction, strong induction, and the well-ordering principle are
logically equivalent over a Peano system. Each can be used as the primitive
proof engine and the others recovered as theorems. This equivalence matters
structurally: it means that induction proofs, strong-induction proofs, and
least-counterexample proofs are different presentations of the same
foundational mechanism.
\end{exposition}

\begin{theorembox}{Theorem (Induction, Strong Induction, and Well-Ordering Are Equivalent)}
\begin{theorem}[Induction, Strong Induction, and Well-Ordering Are Equivalent]
\label{thm:induction-well-ordering-equivalence}
Let $(P,S,1)$ be a Peano system equipped with the order $\le$. The following
are equivalent:
\begin{enumerate}
  \item The induction principle: for every predicate $\varphi$ on $P$, if
  $\varphi(1)$ holds and the successor step holds, then $\varphi(n)$ holds
  for all $n \in P$.
  \item The strong induction principle: for every predicate $\varphi$ on $P$,
  if $\varphi(k)$ holds for every $k<n$ implies $\varphi(n)$, then
  $\varphi(n)$ holds for every $n \in P$.
  \item The well-ordering principle: every nonempty subset of $P$ has a least
  element with respect to $\le$.
\end{enumerate}
\hyperref[prf:induction-well-ordering-equivalence]{Go to proof.}
\end{theorem}
\end{theorembox}

\begin{remark*}[Standard quantified statement]
\[
\begin{aligned}
&\text{Let } (P,S,1) \text{ be a Peano system.}\\
&\operatorname{WellOrderingPrinciple}
\Longleftrightarrow
\operatorname{StrongInduction}
\Longleftrightarrow
\Bigl(\forall \varphi\;
\bigl(\varphi(1) \land
\forall n \in P\;(\varphi(n) \Longrightarrow \varphi(S(n)))
\Longrightarrow
\forall n \in P\;\varphi(n)\bigr)\Bigr).
\end{aligned}
\]
\end{remark*}

\begin{remark*}[Interpretation]
Induction proves well-ordering by contrapositive: if a nonempty set $A$ had
no least element, the set of non-members of $A$ would be inductive, forcing
$A = \varnothing$, a contradiction. Well-ordering proves induction by forming
the set of elements where $\varphi$ fails and, if nonempty, extracting a
least failure to derive a contradiction with the successor step.
\end{remark*}

\begin{dependencies}
\begin{itemize}
  \item \hyperref[thm:induction-principle-for-peano-system]{Induction Principle for a Peano System}
  \item \hyperref[thm:strong-induction-on-peano-system]{Strong Induction on a Peano System}
  \item \hyperref[thm:well-ordering-principle]{Well-Ordering Principle}
\end{itemize}
\end{dependencies}
\end{topicbox}
